\documentclass[11pt]{article}
\usepackage[margin=1in]{geometry}
\usepackage{times}
\usepackage{hyperref}
\hypersetup{colorlinks=true, linkcolor=black, urlcolor=blue, citecolor=black}
\title{A Motor That Speeds Up Under Load Is a Wiring Fault, Not a Discovery: Thane Heins and the Perepiteia Effect}
\author{Flyxion \\ \small Independent Researcher}
\date{}
\begin{document}
\maketitle

\begin{abstract}
Thane Heins, a Canadian inventor, has since the mid-2000s promoted a phenomenon he calls the Perepiteia effect, in which an electric motor-generator, when its output is connected to a load, appears to accelerate rather than experience the expected retarding back-EMF drag, footage of which drew attention from several university engineering departments and brief media coverage suggesting a challenge to Lenz's law. This paper argues that Lenz's law, a direct consequence of energy conservation applied to electromagnetic induction, is not a phenomenological rule of thumb but a theorem following from Faraday's law and the conservation of energy, that the specific acceleration behavior demonstrated in Heins's videos is consistent with a well known category of motor control and wiring configuration artifact in which a motor's own drive circuit responds to load changes in a way that increases rather than decreases input power delivery, producing an apparent but illusory violation, and that no independent university evaluation of the device, despite initial interest, has published a peer-reviewed confirmation of a genuine Lenz's law violation in the years since the claim began circulating.
\end{abstract}

\section{Introduction}

Thane Heins demonstrated, beginning around 2006, a permanent-magnet motor-generator setup in which connecting an electrical load to the generator's output coil appeared, in video demonstrations, to cause the attached drive motor to speed up rather than slow down, the opposite of the retarding effect Lenz's law predicts should occur when a generator delivers current to a load, since the induced magnetic field from that current should oppose the change producing it and thus oppose, rather than assist, the driving motor's rotation. The demonstrations drew attention including interest from engineering researchers at Canadian universities, some of whom conducted preliminary informal examinations, generating media coverage suggesting a possible challenge to a foundational electromagnetic principle.

\section{Why Lenz's Law Is Not an Empirical Rule Awaiting a Counterexample}

Lenz's law, that an induced current's magnetic field opposes the change in flux that produced it, is not an independently posited empirical regularity but a direct theorem derivable from Faraday's law of induction combined with the conservation of energy: were induced currents to instead reinforce rather than oppose the inducing change, a generator connected to a load would spontaneously accelerate its own driving mechanism, extracting unlimited energy from nothing, a runaway process with no energy source, which is precisely why the law's sign convention is fixed by energy conservation rather than by an independent empirical observation that could in principle have come out the other way. A genuine violation of Lenz's law is not merely surprising in the way a new particle or an unexpected material property might be; it would require conservation of energy itself to fail in ordinary electromagnetic induction, a domain confirmed to extraordinary precision across the entire discipline of electrical engineering.

\section{What Motor Control Artifacts Can Produce This Appearance}

Permanent-magnet motor-generator systems using electronic speed controllers or specific brush and commutation configurations can exhibit behavior in which connecting a load changes the electrical characteristics seen by the driving circuit in a way that causes the controller to deliver more input power to the drive motor, for instance if load-induced back-EMF changes alter a feedback signal the controller uses to regulate motor speed, producing a genuine speed increase that is fully accounted for by increased power draw from the motor's own power supply rather than by any anomalous energy contribution from the generator side, a well documented category of artifact in motor control engineering distinct from, but easily visually confused with, a genuine violation of the underlying electromagnetic principle, particularly in demonstration setups that do not include an independent, clearly displayed measurement of total input power to the driving motor throughout the load-connection event.

\section{The Absence of a Peer-Reviewed Confirmation}

Despite initial engineering interest and several universities agreeing to examine early versions of the device, no independent, peer-reviewed publication confirming a genuine Lenz's law violation under properly instrumented conditions, specifically including continuous measurement of total electrical input power to the drive motor throughout the demonstrated acceleration event, has appeared in the years since the claim began circulating, a gap consistent with the motor control artifact explanation and inconsistent with a confirmed genuine anomaly, which would represent a sufficiently significant finding to warrant rapid and thorough independent publication once verified.

\section{The Entropic Gravity Framing}

The Perepiteia effect concerns electromagnetic induction and energy conservation rather than gravitation directly, and is included in this series as a comparative case in the broader "apparent violation of a foundational conservation-adjacent law" genre alongside the EmDrive and Mach effect thruster entries addressed elsewhere here; the specific objection in this case is unusually direct, since Lenz's law's status as a theorem of energy conservation, rather than an independently falsifiable empirical rule, means a genuine violation would carry implications for the conservation of energy itself, a bar considerably higher than the claim's demonstration videos have been shown to clear.

\section{Objections}

Heins and supporters argue that the effect has been observed by multiple independent engineers who examined early prototypes and found the acceleration behavior genuinely puzzling, taking this puzzlement as evidence the effect resists a simple mundane explanation. Initial puzzlement by engineers examining a demonstration video or a prototype not yet fully instrumented for total input power measurement is consistent with, rather than evidence against, an artifact requiring careful instrumentation to identify, and the relevant question is not whether early examiners found the demonstration surprising but whether a subsequent, properly instrumented, peer-reviewed test has confirmed the effect, which none has in the years since the claim's initial circulation.

\section{Conclusion}

Lenz's law follows directly from energy conservation rather than standing as an independently falsifiable empirical rule, the specific acceleration behavior demonstrated is consistent with a well documented motor control artifact category involving increased input power delivery under changing feedback conditions, and no independent peer-reviewed confirmation of a genuine violation has appeared in the years since the claim began circulating, despite the initial engineering interest the demonstration videos generated.

\end{document}
